\section{Product replacement at a cut}

The product replacement argument of~\cite{Schuch2008EntropyScaling}
(arXiv:0705.0292v2, lines 436--470) can be expressed directly in terms of
coefficient spaces and cut ranks.

\begin{lemma}[Product replacement does not increase cut ranks]
    \label{lem:chain_product_replacement_ranks}
    \lean{TNLean.MPS.AreaLaw.rank_cutMatrix_le_of_product_mem_ranges}
    \leanok
    \uses{def:chain_open_mps}
    Let $\psi,\chi$ be vectors on $L$ sites of local dimension $q$, and
    write $C=M_k(\psi)$ for the coefficient matrix at a cut $k$.
    Suppose $u\in\operatorname{range}C$, $v\in\operatorname{range}C^{\mathsf T}$,
    and $M_k(\chi)=uv^{\mathsf T}$. Then
    $\operatorname{rank}M_j(\chi)\le\operatorname{rank}M_j(\psi)$ at every cut $j$.
    No nonvanishing or normalization is required.
\end{lemma}
\begin{proof}\leanok
    \uses{thm:hastings_supported_cut_rank}
    Choose $w,z$ with $Cw=u$ and $C^{\mathsf T}z=v$. The identities
    \begin{align}
        \label{eq:chain_product_replacement_local_maps}
        (uz^{\mathsf T})C=uv^{\mathsf T}
        =C(vw^{\mathsf T})^{\mathsf T}
    \end{align}
    express $\chi$ both as a left-local and as a right-local linear image of
    $\psi$. For $j\le k$, use the right-local image; for $j>k$, use the
    left-local image. The corresponding operator acts on only one side of
    cut $j$, so cannot increase its rank. If there are no configurations,
    both vectors vanish and the assertion is immediate.
\end{proof}

\begin{corollary}[Preservation of the bond bound]
    \label{cor:chain_product_replacement_mps}
    \lean{TNLean.MPS.AreaLaw.IsOpenMPS.of_product_mem_cut_ranges}
    \leanok
    \uses{lem:chain_product_replacement_ranks,def:chain_open_mps}
    Under the hypotheses of Lemma~\ref{lem:chain_product_replacement_ranks},
    if $\psi$ is an open-boundary MPS of bond dimension at most $D$, then
    so is $\chi$. This includes $D=0$ and empty or one-site chains.
\end{corollary}
\begin{proof}\leanok
    \uses{lem:chain_product_replacement_ranks,thm:chain_cut_rank_of_mps}
    Each internal cut rank of $\chi$ is at most the corresponding rank of
    $\psi$, hence at most $D$. Apply the cut-rank characterization.
\end{proof}

\begin{theorem}[Normalized product replacement]
    \label{thm:chain_normalized_product_replacement}
    \lean{TNLean.MPS.AreaLaw.IsOpenMPS.exists_product_overlap_ge}
    \leanok
    \uses{def:chain_open_mps}
    Let $\psi$ be a unit open-boundary MPS of bond dimension at most $D$.
    At any cut $k$, let $a,b$ be arbitrary target vectors on the two sides.
    There are unit vectors $u,v$ and a unit open-boundary MPS $\chi$ of
    bond dimension at most $D$ such that $M_k(\chi)=uv^{\mathsf T}$,
    $\operatorname{rank}M_j(\chi)\le\operatorname{rank}M_j(\psi)$ for every $j$, and
    $|\langle a\otimes b,\psi\rangle|
        \le |\langle a,u\rangle|\,|\langle b,v\rangle|$.
    The inner product on the left is taken in cut coordinates.
\end{theorem}
\begin{proof}\leanok
    \uses{cor:chain_product_replacement_mps}
    The product-overlap theorem in~\cite{QICLean2026Blueprint} gives unit
    $u\in\operatorname{range}M_k(\psi)$ and
    $v\in\operatorname{range}M_k(\psi)^{\mathsf T}$ with the stated overlap
    bound. Transport $u\otimes v$ back to the chain coordinates to define
    $\chi$. This transport preserves norms, so $\|\chi\|=\|u\|\,\|v\|=1$.
    Corollary~\ref{cor:chain_product_replacement_mps} preserves the bond bound.
\end{proof}
